\chapter{Step 2 — Determining Who or What Is Affected}
\label{chap:Step 2 — Determining Who or What Is Affected}

%---------------------------- SECTION ----------------------------
\section{More Than a Formality}
\paragraph{In many risk assessment templates and guidance documents, the question of who or what is affected by an identified risk appears as a single column in a table — a brief note to be filled in and moved past quickly on the way to the more analytically demanding work of evaluation and control. It is often treated, in practice, as little more than a labelling exercise.}
\paragraph{\textbf{This is a mistake — and it is a mistake with real consequences}.}
\paragraph{Thinking carefully and seriously about who or what is exposed to each identified risk, and in what ways, is not a bureaucratic step. It is a substantive analytical exercise that shapes everything that follows. It determines how a risk should be evaluated — because the same hazard presents very differently depending on who is exposed to it and in what circumstances. It determines what controls are appropriate — because effective controls need to be designed around the nature of the exposure, not just the nature of the hazard. And it carries direct regulatory significance — because many of the compliance obligations that matter most to our audience are specifically framed around the protection of defined groups of people or interests.}
\paragraph{This chapter explores who and what should be considered, why the analysis is more nuanced than it first appears, and how getting it right produces a materially better risk assessment.}

%---------------------------- SECTION ----------------------------
\section{Broadening the Frame — Who and What Can Be Affected}
\paragraph{The natural starting point for most people when considering who is affected by a risk is the people most immediately and obviously in the picture — usually employees, or the customers directly involved in a transaction or process. That starting point is reasonable, but it is rarely sufficient.}
\paragraph{A more complete analysis considers the full range of parties and interests that could be affected if a risk materialised — and the full range of things, not just people, that carry value and can be harmed.}

\subsection{People}
\paragraph{The most direct and in many contexts the most significant category of those affected. When considering people, the analysis should extend beyond the most immediately obvious group to include:}
\paragraph{\textbf{Employees and workers} — including not just permanent staff but contractors, agency workers, volunteers, and others who perform work for or on behalf of the organisation. Regulatory requirements around workplace risk typically extend to all workers, not just those on permanent employment contracts, and the distinction between categories of worker has been the subject of significant legal development in recent years.}
\paragraph{\textbf{Customers and clients} — the people who use the organisation's products or services, whether as individual consumers or as representatives of other organisations. In regulated sectors, the protection of customers is frequently the primary purpose of the regulatory framework, and failure to adequately consider customer exposure in a risk assessment is likely to attract regulatory criticism. The FCA's Consumer Duty, discussed in Chapter~\ref{chap:Legal and Regulatory Considerations}, is an explicit regulatory expression of this expectation — firms are required to think actively and rigorously about the outcomes their products and services deliver for customers.}
\paragraph{\textbf{Third parties and members of the public} — people who are not in a direct relationship with the organisation but who could nevertheless be affected by its activities. This might include residents near a manufacturing facility, road users affected by a logistics operation, or members of the public who interact with a regulated firm's financial promotions without becoming customers.}
\paragraph{\textbf{Counterparties and business partners} — other organisations and individuals with whom the organisation has contractual or commercial relationships. Their exposure to risk arising from the organisation's activities — and the organisation's exposure to risk arising from theirs — is a dimension that risk assessments in commercial and financial services contexts in particular need to consider carefully.}
\paragraph{\textbf{Regulators and supervisory bodies} — whilst regulators are not typically thought of as parties who can be harmed by an organisation's risks, the regulatory relationship is affected by the organisation's risk profile. A significant risk event that involves harm to customers, markets, or the financial system will draw regulatory attention and response — and that response is itself a consequential outcome that the organisation needs to consider}
\paragraph{\textbf{Shareholders and investors} — financial losses, regulatory penalties, and reputational damage all ultimately affect those with a financial interest in the organisation. For listed companies and regulated entities with external investors, the impact of risk events on shareholders and investor confidence is a material consideration.}
\paragraph{\textbf{Future generations and wider society} — particularly relevant in the context of environmental, climate, and social risks, where the effects of an organisation's activities may extend well beyond its immediate stakeholders and into the longer term. The growing focus on ESG obligations, discussed in Chapter~\ref{chap:Types of Risk}, reflects a regulatory and societal expectation that organisations consider these broader impacts.}

\subsection{Vulnerable Groups}
\paragraph{Within the category of people affected, \textbf{vulnerable groups} deserve specific and careful consideration. Vulnerability can take many forms — it may be inherent or situational, permanent or temporary — and the regulatory expectation in many sectors is that organisations actively consider the needs and exposures of vulnerable individuals rather than treating all people as equally affected.}
\paragraph{In financial services, the FCA has published detailed guidance on the fair treatment of vulnerable customers, defining vulnerability broadly to include people with health conditions, those experiencing life events such as bereavement or job loss, those with low financial resilience, and those with low capability in areas relevant to the product or service. A risk assessment that fails to consider how identified risks may affect vulnerable customers differently from the general population is likely to be considered inadequate by the FCA.}
\paragraph{In healthcare, education, social care, and other sectors dealing directly with individuals who may be less able to protect their own interests, the consideration of vulnerability is similarly central to the regulatory framework.}
\paragraph{More broadly, the principle that some people are more exposed to certain risks than others — and that this differential exposure should be reflected in the risk assessment — is one that applies across virtually every sector and risk type. Age, disability, language, financial literacy, health status, and situational factors can all affect the degree to which an individual is exposed to a given risk and the nature of the harm they might experience.}

\subsection{Assets and Resources}
\paragraph{Beyond people, a thorough analysis of what is affected should consider the full range of organisational assets and resources that could be harmed if a risk materialised:}
\paragraph{\textbf{Financial assets} — cash, investments, receivables, and other financial resources. Financial harm is often the most immediately quantifiable consequence of a risk event, and understanding which financial assets are exposed is important both for risk evaluation and for understanding the organisation's financial resilience.}
\paragraph{\textbf{Physical assets} — premises, equipment, infrastructure, and inventory. Physical damage, loss, or unavailability of these assets can disrupt operations, generate recovery costs, and in some cases create liability to third parties.}
\paragraph{\textbf{Data and information} — one of the most significant and most rapidly growing categories of organisational asset. Data breaches, data loss, and data corruption can cause harm simultaneously to the individuals whose data is affected, to the organisation's regulatory standing, and to its commercial position. The regulatory framework around data protection, discussed in Chapter~\ref{chap:Legal and Regulatory Considerations}, reflects the specific and serious nature of data as an asset requiring active protection.}
\paragraph{\textbf{Intellectual property} — patents, trademarks, trade secrets, and proprietary processes. Exposure of intellectual property through security failures, contractual breaches, or regulatory non-compliance can have significant long-term commercial consequences.}
\paragraph{\textbf{Technology and systems} — the digital infrastructure and software that organisations increasingly depend on to operate. System failures, cyber attacks, and technology obsolescence can affect not only the organisation's own operations but also the people and organisations that depend on those systems.}
\paragraph{\textbf{Regulatory permissions and licences} — for regulated organisations, the ability to operate depends on maintaining regulatory authorisation. Loss of, or restriction to, a regulatory permission is one of the most serious consequences that can flow from a risk event, and it is an asset — in the broadest sense — that is frequently underrepresented in risk assessment analysis.}

\subsection{Organisational Interests}
\paragraph{Finally, a complete analysis of what is affected should explicitly consider the broader interests of the organisation itself — the things that, whilst not easily categorised as assets, carry real and significant value:}
\paragraph{\textbf{Reputation and public standing} — as explored in Chapter~\ref{chap:Types of Risk} and as we will examine in depth in Chapter~\ref{chap:Reputational Risk Assessment}, reputational damage is a consequence of risk events that can be more enduring and more significant in its practical impact than the immediate financial or regulatory consequences.}
\paragraph{\textbf{Culture and employee morale} — significant risk events, particularly those involving misconduct, regulatory sanction, or major operational failures, can have profound effects on the culture and morale of an organisation's workforce. These effects are difficult to quantify but genuinely consequential for the organisation's longer-term performance and resilience.}
\paragraph{\textbf{Relationships and trust} — with customers, regulators, investors, suppliers, and the wider market. Trust, once damaged, is slow and difficult to rebuild — and many of the most significant practical consequences of risk events flow from the erosion of trust rather than from the immediate harm itself.}

%---------------------------- SECTION ----------------------------
\section{The Concept of Exposure}
\paragraph{Having identified who and what could be affected, the next analytical task is to understand the nature and extent of their \textbf{exposure} — the degree to which they are actually at risk if the hazard materialises.}
\paragraph{Exposure is shaped by several factors:}

\subsection{Proximity}
\paragraph{How close is the potentially affected party to the source of the hazard? Proximity can be physical — people working directly with a hazardous process are more exposed than those working nearby — but in a compliance and legal context, it is more often conceptual. A customer who is the direct recipient of a financial product is more exposed to conduct risk than one who is a secondary beneficiary. An employee with direct access to sensitive data is more exposed to insider threat risk than one without.}

\subsection{Frequency and Duration}
\paragraph{How often and for how long is the potentially affected party exposed to the hazard? A risk that affects a large number of people briefly may have a very different overall impact profile to one that affects a smaller number of people over an extended period. Frequency and duration are important inputs into the evaluation of likelihood and consequence that follows in the next step.}

\subsection{Capacity to Protect}
\paragraph{To what extent can the potentially affected party protect themselves from the risk — and to what extent does the organisation have a responsibility to do so on their behalf? A sophisticated institutional investor may be well placed to assess and manage the risks of a complex financial product. A retail consumer with limited financial literacy may not be. The degree to which the affected party can protect themselves shapes both the ethical and the regulatory dimension of the organisation's obligations.}
\paragraph{This consideration connects directly to the principle of vulnerability discussed above — and to the broader regulatory expectation that organisations do not simply assume that their customers, employees, or other stakeholders can look after themselves}

\subsection{Interdependencies}
\paragraph{In practice, the exposure of one party is often connected to the exposure of others. A failure in one part of a supply chain creates exposure for all parties downstream. A data breach affecting one organisation may expose the customers of multiple connected organisations. A regulatory failure by one firm in a market may affect the reputation of the entire sector.}
\paragraph{Understanding these interdependencies — the ways in which exposure is transmitted and amplified across networks of relationships — is increasingly important as organisations become more interconnected and as regulators focus more attention on systemic risk.}

%---------------------------- SECTION ----------------------------
\section{Mapping Affected Parties to Identified Risks}
\paragraph{In practice, not every identified risk will affect every potential stakeholder. Part of the analytical work in this step is to map each identified risk to the specific parties and interests that are most directly and significantly exposed — creating a clear picture of who is at risk from what.}
\paragraph{This mapping exercise serves several purposes:}
\paragraph{\textbf{It sharpens the subsequent evaluation} — understanding who is exposed, and how, provides crucial context for assessing the significance of a risk. A risk that primarily affects the organisation's own financial position may be evaluated very differently from one that primarily affects the wellbeing of vulnerable customers.}
\paragraph{\textbf{It identifies regulatory obligations} — many regulatory requirements are triggered by the nature of the affected party. Requirements around consumer protection, data subject rights, employee safety, and environmental harm all flow from the identification of specific categories of affected party.}
\paragraph{\textbf{It informs control design} — controls need to be designed around the nature of the exposure, not just the nature of the hazard. A control that effectively protects the organisation's financial assets may do nothing to protect the customers who are simultaneously exposed to the same underlying risk.}
\paragraph{\textbf{It supports proportionate prioritisation} — risks that affect large numbers of people, or that disproportionately affect vulnerable groups, generally warrant higher priority in the evaluation and control steps that follow.}
\paragraph{A useful format for capturing this mapping is to add to each risk entry in the developing risk register:}
\begin{itemize}
    \item{A list of the primary affected parties — the people, assets, and interests most directly and significantly exposed.}
    \item{A brief description of the nature of their exposure — what specifically could happen to them if the risk materialised.}
    \item{Any particular vulnerability or capacity considerations relevant to the affected parties.}
    \item{Any regulatory obligations triggered by the nature of the affected parties.}
\end{itemize}

%---------------------------- SECTION ----------------------------
\section{The Regulatory Dimension — Obligations Shaped by Who Is Affected}
\paragraph{For compliance and legal professionals, one of the most practically significant aspects of this step is its direct connection to regulatory obligations. Many of the most important compliance requirements are not framed in terms of the risk type — they are framed in terms of the people or interests that the regulation is designed to protect.}
\paragraph{Consumer protection requirements exist because customers are affected by firm behaviour. Data protection requirements exist because individuals are affected by the processing of their personal data. Employment law requirements exist because workers are affected by the decisions and conditions their employers create. Environmental requirements exist because communities, ecosystems, and future generations are affected by industrial and commercial activity.}
\paragraph{This means that the process of identifying who is affected by each risk is not merely a descriptive exercise — it is a direct input into the identification of applicable regulatory obligations. A risk that affects personal data triggers GDPR obligations. A risk that affects retail consumers in a regulated market triggers conduct of business obligations. A risk that affects workers triggers health and safety obligations. Getting this connection right is part of what it means to conduct a risk assessment that is genuinely responsive to the regulatory environment.}

%---------------------------- SECTION ----------------------------
\section{A Practical Illustration — Walking Through an Example}
\paragraph{To bring this step to life, consider the following scenario. A mid-sized financial services firm is introducing a new automated decision-making system to process customer credit applications. The risk identification step has produced a list of risks associated with this system. Now, at Step 2, the firm needs to consider who and what is affected by each of those risks.}
\paragraph{Take one specific identified risk: \textbf{the possibility that the automated system produces biased or discriminatory outcomes, rejecting applications from certain groups at a higher rate than others}.}
\paragraph{Working through the affected parties analysis:}
\paragraph{\textbf{Customers} — most directly affected. Those whose applications are incorrectly rejected suffer real and concrete harm — denial of access to credit, potential financial distress, and the indignity of discriminatory treatment. Within the customer group, the analysis should consider whether particular sub-groups — defined by protected characteristics such as age, race, disability, or gender — are disproportionately exposed.}
\paragraph{\textbf{Vulnerable customers} — a specific and significant sub-group. Individuals who are already in financial difficulty, who have limited access to alternative credit sources, or who lack the capacity to challenge a decision effectively are more severely harmed by an incorrect rejection than those with more resources and resilience.}
\paragraph{\textbf{The organisation itself} — exposed to regulatory, legal, reputational, and financial consequences. Regulatory obligations under equality legislation, consumer protection rules, and data protection law (automated decision-making is specifically addressed in the GDPR) are all engaged. Litigation risk from affected customers is real. Reputational damage from adverse media coverage or regulatory action is a further exposure.}
\paragraph{\textbf{The regulator} — the FCA's objective of ensuring that consumers receive fair outcomes means that discriminatory automated decision-making is a direct concern for the regulator. The firm's regulatory relationship is at risk if this issue materialises and is seen to reflect inadequate governance.}
\paragraph{\textbf{The broader market} — if discriminatory automated decision-making becomes a sector-wide issue, the reputational and regulatory consequences extend beyond any individual firm.}
\paragraph{This analysis immediately reveals that the risk is more complex and more significant than a simple description of the hazard would suggest. It triggers specific regulatory obligations — around equality, data protection, and consumer protection. It identifies vulnerable groups deserving particular consideration. And it reveals a range of affected interests beyond the immediate customer relationship.}
\paragraph{That is the value of doing this step properly.}

%---------------------------- SECTION ----------------------------
\section{Common Shortcomings in Practice}
\paragraph{As with hazard identification, there are characteristic patterns of failure in the analysis of who and what is affected. Being aware of them helps to avoid them.}
\paragraph{\textbf{Focusing exclusively on the organisation itself} — many risk assessments, particularly those conducted by smaller or less mature organisations, focus primarily on the consequences for the organisation — financial loss, regulatory penalty, reputational damage — without adequately considering the harm to external parties. This inward focus is understandable but is likely to produce both an incomplete risk assessment and one that does not adequately reflect the organisation's regulatory obligations.}
\paragraph{\textbf{Treating all affected parties as equivalent} — not all affected parties are equally exposed, equally vulnerable, or equally able to protect themselves. A risk assessment that lists affected parties without differentiating between them misses important analytical nuance that should inform evaluation and control.}
\paragraph{\textbf{Overlooking indirect and systemic effects} — the immediate and direct effects of a risk event are usually the easiest to identify. The indirect effects — the downstream consequences for parties not immediately involved, the systemic effects of a failure that ripples through a network of relationships — are harder to see but can be equally or more significant.}
\paragraph{\textbf{Conflating the affected party with the risk owner} — the person or team responsible for managing a risk is not necessarily the person most affected by it. Conflating these two concepts — confusing ownership with exposure — can lead to control designs that protect the risk owner rather than those who are genuinely at risk.}
\paragraph{\textbf{Neglecting non-human affected interests} — the environment, ecosystems, communities, and future generations are increasingly recognised as legitimate interests that risk assessments should consider. Organisations that fail to include these interests in their analysis risk both ethical and regulatory shortcomings as ESG expectations continue to develop.}

%---------------------------- SECTION ----------------------------
%\newpage
\section{Chapter Summary}
In this chapter we learned about the following key points:
\begin{table}[H]
  \centering
  \renewcommand{\arraystretch}{1.5}
  \begin{tabularx}{\textwidth}{ | >{\raggedright\arraybackslash}p{0.35\textwidth} 
                                | >{\raggedright\arraybackslash}X | }
    \hline
    \textbf{Category of Affected Party} & \textbf{Examples} \\
    \hline
    \textbf{Employees and workers} & All categories of worker, including contractors and agency staff\\
    \hline
    \textbf{Customers and clients} & Direct recipients of products and services, including vulnerable groups\\
    \hline
    \textbf{Third parties and public} & Those affected without a direct relationship with the organisation\\
    \hline
    \textbf{Counterparties and partners} & Other organisations in commercial or contractual relationships\\
    \hline
    \textbf{Shareholders and investors} & Those with financial interests in the organisation\\
    \hline
    \textbf{Wider society} & Communities, future generations, and the public interest\\
    \hline
    \textbf{Financial assets} & Cash, investments, and other financial resources\\
    \hline
    \textbf{Data and information} & Personal data, proprietary information, and intellectual property\\
    \hline
    \textbf{Technology and systems} & Digital infrastructure and operational systems\\
    \hline
    \textbf{Regulatory permissions} & Licences and authorisations required to operate\\
    \hline
    \textbf{Reputation and trust} & Organisational standing with key stakeholders\\
    \hline
  \end{tabularx}
  \caption{Affected Parties}
\end{table}

\begin{table}[H]
  \centering
  \renewcommand{\arraystretch}{1.5}
  \begin{tabularx}{\textwidth}{ | >{\raggedright\arraybackslash}p{0.35\textwidth} 
                                | >{\raggedright\arraybackslash}X | }
    \hline
    \textbf{Factor} & \textbf{Why It Matters} \\
    \hline
    \textbf{Proximity} & Closeness to the hazard source shapes the nature and degree of risk\\
    \hline
    \textbf{Frequency and duration} & How often and how long exposure occurs affects overall impact\\
    \hline
    \textbf{Capacity to protect} & Ability to self-protect shapes both ethical and regulatory obligations\\
    \hline
    \textbf{Interdependencies} & Exposure is often transmitted and amplified across connected relationships\\
    \hline
  \end{tabularx}
  \caption{Factors Shaping Exposure}
\end{table}
\newpage
\paragraph{Key takeaways:}
\begin{itemize}
  \item{Determining who and what is affected is a substantive analytical step — not a labelling exercise.}
  \item{The analysis should extend well beyond the most immediately obvious affected parties to include indirect, systemic, and broader societal effects.}
  \item{Vulnerable groups deserve specific and deliberate consideration — regulatory expectations in many sectors make this explicit.}
  \item{The nature of the affected parties directly determines which regulatory obligations are engaged — making this step a critical input into compliance analysis.}
  \item{Exposure varies significantly between affected parties — proximity, frequency, capacity, and interdependency all shape the picture.}
  \item{The output of this step enriches the risk register and directly informs the evaluation and control work that follows.}
\end{itemize}

\barquote{In Chapter~\ref{chap:Step 3 — Evaluating and Analysing Risks}, we move into the analytical heart of the process — evaluating and analysing the risks we have identified and mapped, developing a clear and defensible picture of their likelihood and potential consequences, and building the prioritised risk profile that will drive the decisions that follow.}